\scanpage{11}
\[
\rho(n)\overset{\mathrm{dfn}}{=}\operatorname{rg}_k L_n
=\sum_{1\leq i\leq n+1}\mu_i
\underset{\substack{\downarrow\\\text{pol. de degré}\\i\ \text{en }n}}{\binom{n+1}{i}}
+\sum_{0\leq j\leq n}\nu_j
\underset{\substack{\downarrow\\\text{pol. de degré}\\j\ \text{en }n}}{\binom n j}
\]
où $\mu_i=\operatorname{rg}_k M_i$, $\nu_j=\operatorname{rg}_k N_j$

donc si \deleted{$\mu_i$ (resp $\nu_j$) \unsure{$=$}} un $\mu_i$ ou un $\nu_j$ est $\ne0$, cela montre que $\rho(n)$ pour $n$ grand est minoré par un polynôme de degré $i$ resp $j$.

[2°) Cas de $\Lambda^k$. Utilisant la formule de $\Lambda^k$ d'une somme directe, et 1°), on est ramené à prouver que les $\Lambda^k$ d'un module $\frac12$ simplicial $L_{\bullet}$ de la forme $\Lambda^i\Phi_*\otimes M_{\unsure{i|{}}}$, $\Lambda^j\Psi_*\otimes N$. Comme $M,N$ est \unsure{facteur|facteurs} direct d'un libre $P$, $L_{\bullet}$ est facteur direct de $L'_{\bullet}=\Lambda^i\Phi_*\otimes P$ ou $\Lambda^j\Psi_*\otimes P$, \deleted{\unsure{et que}} \deleted{\unsure{prop. 1 ci-dessus est}} et en réutilisant la formule de $\Lambda^k$ d'une somme directe, \unsure{montrant|montre} que $\unsure{\Gamma|\Lambda}^k(L_{\bullet})$ est facteur direct de $\Lambda^k L'_{\bullet}$. \inserted{Comme la propriété du \unsure{lemme} \deleted{\unsure{de}} de \unsure{prop.} 1, son cor. 1 et du cor 2 \deleted{sont} stable par fact directs,} on est ramené à prouver $\Lambda^i L'_{\bullet}$ est \deleted{\unsure{str}} spécial, donc \unsure{sps|supp.} $M,N$ libres --- donc $\overline{\mathrm m}$ (par la formule de la somme) $M=k$, $N=k$. Donc on est ramené à prouver
% archive sheet number: 10
